\subsection{(\S1) The Eight Function-Attitudes Unpacked}

\begin{enumerate}

    \item \ovalbox{\it MBTI and function-attitudes} To help people apply their
        scores on the MBTI® to themselves, and to reap the deeper
        transformative benefits of Jungian type theory, we need to be able to
        recognize the different types of consciousness Jung originally
        described. Practically, that means we have to be prepared to recognize
        in someone's interactive or introspective personality style introverted
        feeling (Fi), introverted thinking (Ti), introverted intuition (Ni),
        introverted sensation (Si), extraverted feeling (Fe), extraverted
        thinking (Te), extraverted intuition (Ne), or extraverted sensation
        (Se), simply on the basis of what the person reveals in our presence.
        We also need to be able to assess how well the function-attitude
        displayed is actually working for the person who is showing it to us.

    \item \ovalbox{\it type understanding is challenging} My own experience as a
        clinician who believes that type understanding is crucial to uncovering
        what limits people and what helps them to develop is that very few us,
        whether clients or type practitioners, can actually recognize the eight
        function-attitudes.

    \item \ovalbox{\it novelist challenge to Beebe} The eight function-attitudes
        that he describes in Psychological Types (1921/1971, ¶¶556–671) were
        for him the 'psychological types.' This term meant something quite
        specific for him: the eight types of awareness capable of constituting
        a psychological individual's conscious psychology. They are mental
        processes that, even if still in the unconscious, are likely to press
        for integration into the psyche, where they can become, for the first
        time, consciousnesses. They hold the motivation to energize will and to
        develop over time—will and time being additional factors that can
        become linked to consciousness.

\end{enumerate}

\clearpage

\begin{table}[ht]
    \caption[Semantic Field of Each Function-Attitude.~\verified]{Semantic Field of Each Function-Attitude. The reader may recall that each of these words is called in grammar a gerund, a verbal noun or process word. Each describes a procedure that a mind can follow over time, a procedure aimed at perceiving or judging reality as accurately as possible, which is the aim of any consciousness.\\
Definitions: e = Extraverted, i = Introverted; S = Sensation, N = Intuition, T = Thinking, F = Feeling.\\
\verified~\textcolor{gray}{\small Source: Page 4-7, Section ``The Eight Function-Attitudes Unpacked'' from Beebe~\cite{beebe}.}}
    \label{tab:beebe-eight-function}
    \begin{center}
	    \begin{tabular}{cccc}
		    \toprule
		    \multirow{2}{*}{\textbf{Function-attitude}} & \multicolumn{3}{c}{\textbf{Semantic Field (Triangle)}} \tabularnewline
		    \cmidrule(lr){2-4}
		    & \small{Persona} & \small{Ego} & \small{Self} \tabularnewline
		    \midrule

		    \rowcolor{black!10}Se & Engaging & Experiencing & Enjoying \tabularnewline
		    Si & Implementing & Verifying & Accounting \tabularnewline

		    \rowcolor{black!10}Ne & Entertaining & Envisioning & Enabling \tabularnewline
		    Ni & Imagining & Knowing & Divining \tabularnewline

		    \rowcolor{black!10}Te & Regulating & Planning & Enforcing \tabularnewline
		    Ti & Naming & Defining & Understanding \tabularnewline

		    \rowcolor{black!10}Fe & Validating & Affirming & Relating \tabularnewline
		    Fi & Judging & Appraising & Establishing the Value \tabularnewline

		    \bottomrule
	    \end{tabular}
    \end{center}

\end{table}

\begin{description}

\item[Semantic Field] The middle word is the apex of a triangle.  The first of
the three words (as we read across for each of the eight mental processes) can
be thought of as what the process looks like on the surface, to another person
seeing someone for the first time using that process -- what we might call its
\textit{persona} level.  The second, central, word captures the heart of the
process -- the one that is embraced by the \textit{ego} as its chief concern.
The third word embodies what the process becomes at its most evolved level when
it is working in sympathetic harmony with what Jung call the \textit{Self}: one
might call this the goal of the process.

\item[Example] Introverted feeling (Fi) can develop from an initial, near
instinctual `judging' into a mature, reflective `appraising' before evolving to
a more far-seeing `establishing the value.' In this latter stage, which aspires
to wisdom, the feeling reaction finds its ground in a more ideal, archetypal
realm, which allows it to discriminate the value that has led to the earlier
judgments and appraisals.

\item[Suggestion to Learners] To focus on the central words in my horizontal
sequences for the different functions of consciousness, using the wing words as
the beginning and the end of the process whose functioning they are trying to
become conscious of.  Using this method of observation, it will become clearer
the degree to which some people are centrally engaged in appraising, defining,
knowing, or verifying as \ul{the internal process by which they hope to master
the objects they encounter} (or say, internal methods by which they have
already learned to deal with what they encounter).  Other people are more
occupied with affirming, planning, envisioning, or experiencing as \ul{their
way of expressing their willingness to subject themselves to the influence of
the objects with which they are trying to connect} (or say, methods by which
they enact a genuine encounter with the Other).  Seeing this fundamental
difference between the introverted and extraverted functions, we can gradually
come to understand not just the ways the functions differ in their introverted
and extraverted aspects, but what the different purposes are that make them
functions.

\end{description}
